\documentclass[10pt,a4paper]{amsart}
\usepackage[margin=19mm]{geometry}
\usepackage{amsmath,amssymb,mathtools}
\usepackage[scr=rsfs,cal=euler]{mathalfa}
\usepackage{xfrac}
\usepackage[unicode,hidelinks,bookmarks=false]{hyperref}
\newcommand{\ie}{i.e.\ }
\newcommand{\cf}{cf.\ }
\newcommand{\resp}{respectively\ }
\newcommand{\Subsection}{\subsection}
\providecommand{\nobreakdash}{\nobreak\hbox{-}}
\setlength{\emergencystretch}{2em}
\multlinegap0pt
\usepackage{xcolor}
\usepackage{amssymb}
\usepackage[shortlabels]{enumitem}
\usepackage[normalem]{ulem}
\usepackage{tikz}
\newtheorem{theorem}{Theorem}
\newtheorem{claim}{Claim}
\title{Shadowing phenomenon for composition operators on the Hardy space $H^2(\mathbb{D})$}
\author{Artur Blois, Ben-Hur Eidt, Paulo Lupatini, Osmar R. Severiano}
\date{Source: arXiv:2603.10575v2 (CC BY 4.0)}
\begin{document}
\maketitle

\noindent\emph{Source and licence.} Shadowing phenomenon for composition operators on the Hardy space $H^2(\mathbb{D})$. Version arXiv:2603.10575v2 (CC BY 4.0), distributed by arXiv under the Creative Commons Attribution 4.0 International licence, http://creativecommons.org/licenses/by/4.0/, as recorded in the arXiv licence metadata for this exact version and shown on the abstract page https://arxiv.org/abs/2603.10575v2. Full licence notice: \texttt{licenses/arxiv-2603.10575-CC-BY-4.0.txt}.\par\smallskip

\noindent\textit{Editorial scope.} Counted unit: the article's theorem Difícil (source0.tex lines 705--706). The article's complete original proof of it is delivered with it (source0.tex lines 708--796, one \texttt{proof} environment of 89 lines). No mathematics is written, repaired or re-derived: the statement and the proof are the article's own lines, in the article's own notation, and no retained line is substituted or re-flowed.  No further source-local context is needed: every cross-reference of every retained line resolves inside the two delivered spans.  Nothing was omitted from any retained span, so the package carries no omission record.  No retained line contains a citation, so the wrapper carries no bibliography and no bibliography entry.  No retained line contains a figure environment, an image-inclusion command or drawing code, so no image is delivered and the package declares no asset dependency.  Editorial formatting only, in addition: an AMS \texttt{amsart} preamble that restates the article's own declarations and macros used by the retained lines, namely the environments \texttt{theorem}, \texttt{claim} on separate counters, together with the standard packages those lines need.  The retained environments print in this document as Theorem 1, Claim 1 in order of appearance, whereas the article prints the same items with the numbering of its own full document.  The complete native source of the article as distributed in the arXiv e-print for this exact version is bundled unchanged as \texttt{sources/196146/source0.tex}.
\begin{theorem}\label{Teo Difícil} If $a\in (0,1)$ then $W_a$ is a generalized hyperbolic operator.
\end{theorem}

\begin{proof}  Observe that $L^2$ is the direct sum of the following closed subspaces $M=\{F\in L^2:F \ \text{vanishes a.e. on}\ (0,a) \}$ and $N=\{F\in L^2:F\ \text{vanishes a.e. on}\ [a, \infty) \}.$  We will show that $M$ and $N$ satisfy the conditions of the definition of generalized hyperbolicity. 

\begin{claim}
$W_a(M)\subset M$ and $r({W_a}_{|M})<1.$ 
\end{claim}
It is easy to see that $W_a(M)\subset M.$ Inductively, we see that the $n$-th iterate of $W_a$ is given by
\begin{align*}
(W_a^nF)(t)=e^{-t(1-a^n)}F(a^nt),
\end{align*}
and by the change of variable $s=a^nt$ we obtain the following estimate
\begin{align*}
\|W_a^n F\|^2_{L^2} & = \int_{\mathbb{R}_+} |(W_a^nF)(t)|^2 dt =\int_{\mathbb{R}_+} |e^{-t(1 - a^n)} F(a^nt)|^2 dt\\
&=\frac{1}{a^n} \int_{\mathbb{R}_+} e^{-2s( \frac{1}{a^n} - 1)} |F(s)|^2 ds  =  \frac{1}{a^n} \int_{a}^{\infty} e^{-2s( \frac{1}{a^n} - 1)} |F(s)|^2 ds \\
&\leq  \frac{e^{-2a( \frac{1}{a^n} - 1)}}{a^n}  \|F\|_{L^2}^2
\end{align*}
for all $F\in M.$ Hence, for all $n \in \mathbb{N}$
\begin{align*}
\|(W_a)^n_{|M}\|^{\frac{1}{n}} \leq  \frac{e^{-\frac{a}{n}( \frac{1}{a^n} - 1)}}{\sqrt{a}} .
\end{align*}
 Since $a\in(0,1),$ it follows that $\frac{a}{n}( \frac{1}{a^n} - 1) =  \frac{ 1 - a^n}{na^{n - 1}} \to  \infty$ as $n \to \infty.$ Therefore, $r({W_a}_{|M}) = 0$.


\begin{claim} ${W_a}|_{N}: N \to W_a(N)$ is a bijection, $W_a(N)$ is closed,  $N \subset W_a(N)$ and $r((W_a{_{|N}})^{-1})_{|N})< 1.$ 
\end{claim}
%\textcolor{magenta}{Paulo: Eu acho que deveria definir \(E = \{F\in L^2:F\ \text{vanishes a.e. on}\ [1, \infty) \}\) e  escrever \(W_a(N) \subset E\) e depois a mesma coisa pra outra inclusao.}

Since $W_a$ is injective, it follows that ${W_a}_{|N}$ is injective and hence ${W_a}|_{N}: N \to W_a(N)$ is a bijection. In order, let us show  that $W_a(N)$ is closed. We will do this by showing that $W_a(N)=\{F\in L^2:F\ \text{vanishes a.e. on}\ [1, \infty) \}.$ Indeed,  for $F\in N$ we have $F(at)=0$ for a.e. $t\geq1$ because $at \geq a$. Thus $W_aF$ vanishes on $[1,\infty).$ This gives the inclusion $``\subset".$ On the other hand, given $F\in L^2$ that vanishes a.e. on $[1,\infty),$ we define the following function
\begin{align}\label{surjective}
G(t)=\begin{cases}
e^{\frac{1-a}{a}t}F(t/a), & 0<t<a,\\
0,& t\geq a.
\end{cases}
\end{align}
By the change of variable $s=t/a,$ we get
\begin{align*}
\int_{\mathbb{R}_+}|G(t)|^2dt&=\int_{0}^{a} e^{\frac{2(1-a)}{a}t}|F(t/a)|^2dt = a\int^1_0 e^{2s(1-a)}|F(s)|^2ds\\
&\leq ae^{2(1-a)}\int^1_0|F(s)|^2ds\leq ae^{2(1-a)}\|F\|_{L^2}^2
\end{align*}
which shows that $G\in L^2.$ Moreover, we have $G\in N$ and
\begin{align*}
(W_aG)(t)&=e^{-t(1-a)}G(at)=\begin{cases}
e^{-t(1-a)}e^{\frac{1-a}{a}at}F(t), & 0<t<1,\\
0,& t\geq 1
\end{cases}\\
&=F(t).
\end{align*}
Since $W_aG=F,$ we obtain the inclusion $``\supset ".$ In particular, it follows that $N\subset W_a(N).$ Now, for each $F\in N$ we define
\begin{align*}
(V_aF)(t)=e^{\frac{1-a}{a}t}F(t/a).
\end{align*}

%\textcolor{magenta}{Paulo: Eu mudaria um pouco o final da prova: começando aqui(NOTE O \(a^2\), note que tem uma igualdade que antes era desigualdade)}

%Now, for each $F\in N$, we define $(V_aF)(t)=e^{\frac{1-a}{a}t}F(t/a)$. Since $F$ vanishes a.e. on $[a,\infty)$, the function $V_aF$ vanishes a.e. on $[a^2,\infty)$ and hence belongs to $N$. By the change of variable $s=t/a$, we get
%\begin{align*}
  %  \|V_aF\|_{L^2}^2 & =\int_{\mathbb{R}_+}e^{\frac{2t}{a}(1-a)}|F(t/a)|^2dt = a\int_{\mathbb{R}_+}e^{2s(1-a)}|F(s)|^2ds\\ 
 %                    & = a\int_0^a e^{2s(1-a)}|F(s)|^2ds \leq ae^{2a(1-a)}\|F\|_{L^2}^2,
%\end{align*}
%which shows that $V_a:N\to N$ is a well-defined bounded operator. Moreover, a direct computation gives $(W_a|_N)\circ V_a=I_N$. Since $N\subset W_a(N)$, it follows that $V_a=\left.(W_a|_N)^{-1}\right|_N$.

%For every $n\in\mathbb N$, the $n$-th iterate of $V_a$ is given by $(V_a^nF)(t)=e^{\frac{t}{a^n}(1-a^n)}F(t/a^n)$. Using the change of variable $s=t/a^n$, we obtain
%\begin{align*}
%    \|V_a^nF\|_{L^2}^2 & = \int_{\mathbb{R}_+}e^{\frac{2t}{a^n}(1-a^n)}|F(t/a^n)|^2dt\\
%                       & = a^n\int_{\mathbb{R}_+}e^{2s(1-a^n)}|F(s)|^2ds = a^n\int_0^ae^{2s(1-a^n)}|F(s)|^2ds.
%\end{align*}
%Since the essential supremum of $e^{2s(1-a^n)}$ on $(0,a)$ is $e^{2a(1-a^n)}$, it follows that $\|V_a^n\|=a^{n/2}e^{a(1-a^n)}$. Therefore, $r(V_a)=\lim_{n\to\infty}\|V_a^n\|^{1/n}=\sqrt{a}<1$. This completes the proof.

%\textcolor{magenta}{Paulo: até aqui. Abaixo esta a original}
Note that $V_aF$ vanishes a.e. on $[a,\infty)$ (indeed, it vanishes in $[a^2, \infty) \supset [a, \infty)$) and  by the change of variable $s=t/a$ we get  
\begin{align*}
\int_{\mathbb{R}_+}|(V_aF)(t)|^2dt&=\int_{\mathbb{R}_+}e^{\frac{2t}{a}(1-a)}|F(t/a)|^2dt=a\int_{\mathbb{R}_+}e^{2s(1 - a)}|F(s)|^2ds\\
&= a\int_{0}^{a} e^{2s(1-a)}|F(s)|^2ds\leq ae^{2a(1-a)}\|F\|^2_{L^2},
\end{align*}
which shows that $V_a: N \to N$ is a well defined operator. Moreover, a simple computation gives ${W_a}_{|N} \circ V_a=I$  and hence 
\begin{align*}
({W_a}_{|N})^{-1}F=({W_a}_{|N})^{-1}\circ({W_a}_{|N}) \circ V_aF = V_aF
\end{align*}
for each $F\in N.$ This allows us to conclude that $\left( ({W_a}_{|N})^{-1} \right)_{|N}=V_a.$ From this operator equality, we compute the iterates of $({W_a}_{|N})^{-1}_{|N},$ namely:
\begin{align*}
(V_a^nF)(t)=e^{\frac{t}{a^n}(1-a^n)}F(t/a^n).
\end{align*}
We now estimate $\|V_a^n\|.$ By using the change of variable $s=t/a^n$, we obtain
\begin{align*}
\|V_a^nF\|^2_{L^2}&=\int_{\mathbb{R}_+}|(V_a^nF)(t)|^2dt=\int_{\mathbb{R}_+}e^{\frac{2t}{a^n}(1-a^n)}|F(t/a^n)|^2dt\\
&=a^n\int_{\mathbb{R}_+}e^{2s(1-a^n)}|F(s)|^2ds=a^n\int_{0}^ae^{2s(1-a^n)}|F(s)|^2ds\\
&\leq a^ne^{2a(1-a^n)}\|F\|^2_{L^2}
\end{align*}
for all $F\in N,$ which implies $\|V_a^n\|\leq a^{\frac{n}{2}}e^{a(1-a^n)}.$ Therefore, $r(V_a) \leq \sqrt{a}<1.$ This completes the proof.
\end{proof}

\end{document}
